% Solutions/hints for ch13 exercises (assembled into Appendix C)

\begin{solution}{exr:ch13-mirror}
From $B$'s point of view the relative position is $\pos_A - \pos_B = -\pos_{\mathrm{rel}}$ and the relative velocity is $\vel_B - \vel_A = -\vel_{\mathrm{rel}}$, so $\VO^{\ttc}_{B|A} = \set{\vel : \exists t \in (0,\ttc],\ t\vel \in \disc{-\pos_{\mathrm{rel}}}{R}} = -\VO^{\ttc}_{A|B}$. Point reflection through the origin maps the boundary to the boundary and preserves distances, so the boundary point of $\VO^{\ttc}_{B|A}$ closest to $-\vel_{\mathrm{rel}}$ is $-\vect{q}$, the change is $-\vect{q} - (-\vel_{\mathrm{rel}}) = -\vect{u}$, and the outward normal at $-\vect{q}$ is $-\vect{n}$. \Cref{def:ch13-orca} then gives the point $\vel_B + \tfrac12(-\vect{u})$ and the normal $-\vect{n}$. Adding $(\vel_A^{\mathrm{new}} - \vel_A - \tfrac12\vect{u})\cdot\vect{n} \ge 0$ and $(\vel_B^{\mathrm{new}} - \vel_B + \tfrac12\vect{u})\cdot(-\vect{n}) \ge 0$ gives $(\vel_A^{\mathrm{new}} - \vel_B^{\mathrm{new}} - \vel_{\mathrm{rel}} - \vect{u})\cdot\vect{n} = (\vel^{\mathrm{new}}_{\mathrm{rel}} - \vect{q})\cdot\vect{n} \ge 0$, the closed half-plane on the outer side of the tangent line at $\vect{q}$; \cref{thm:ch13-convex} finishes the argument.
\end{solution}

\begin{solution}{exr:ch13-by-hand}
Reflecting the instance in the $x$-axis maps $\pos_{\mathrm{rel}}$ to $(4,-0.5)$ and leaves $\vel_{\mathrm{rel}} = (2,0)$ unchanged, so every number of \cref{tab:ch13-example} is reflected too. Now $\varphi = -7.13^\circ$, $\vect{w} = (1, 0.125)$ and $\pos_{\mathrm{rel}} \times \vect{w} = +1 > 0$, so the closest boundary point is on the \emph{left} leg, with direction $\vect{d}_{\mathrm{L}} = (0.9920, 0.1260)$, foot $\vect{q} = (1.9682, 0.2500)$, $\vect{u} = (-0.0318, 0.2500)$ and outward normal $\vect{n} = (-0.1260, 0.9920)$ (the left-leg normal is $\vect{d}$ rotated by $+90^\circ$). A's half-plane has point $\vel_A + \tfrac12\vect{u} = (0.9841, 0.1250)$ and normal $\vect{n}$, and the program returns $\vel_A^{\mathrm{new}} = (1.1809, 0.1500)$: A passes above B instead of below. The mirroring is exact because the construction uses only distances, dot products and the sign of one cross product, all of which commute with a reflection.
\end{solution}

\begin{solution}{exr:ch13-arc-case}
(a) The ray $t \mapsto t\vel_{\mathrm{rel}}$ with $\vel_{\mathrm{rel}} = (0.4, 0.05)$ points at $7.13^\circ$, almost exactly along $\pos_{\mathrm{rel}}$, so it enters the disc $\disc{\pos_{\mathrm{rel}}}{1}$ when $t\norm{\vel_{\mathrm{rel}}} \approx \norm{\pos_{\mathrm{rel}}} - 1$, that is, at $t \approx 3.03/0.403 \approx 7.5$\,s $> \ttc$: outside $\VO^{\ttc}_{A|B}$ (\code{time\_to\_collision} gives $7.52$\,s). (b) $\vect{c} = (1, 0.125)$, $\rho = 0.25$, $\vect{w} = (-0.6, -0.075)$, $\norm{\vect{w}} = 0.6047$; $\vect{w}\cdot\pos_{\mathrm{rel}} = -2.4375 < 0$ and $(\vect{w}\cdot\pos_{\mathrm{rel}})^2 = 5.941 > R^2\norm{\vect{w}}^2 = 0.366$, so the arc case applies: $\vect{n} = \vect{w}/\norm{\vect{w}} = (-0.9923, -0.1240)$, $\vect{u} = (\rho - \norm{\vect{w}})\vect{n} = (0.3519, 0.0440)$ and $\vect{q} = \vel_{\mathrm{rel}} + \vect{u} = (0.7519, 0.0940)$, the point of the arc nearest to the apex. $\vect{u}\cdot\vect{n} = -0.3547 < 0$: the relative velocity is outside the truncated cone by $0.355$\,m/s, and A's half-plane through $\vel_A + \tfrac12\vect{u}$ allows A to move up to $0.177$\,m/s towards the cone. (c) Let $T$ be the tangent point of the leg with direction $\vect{d}$, so $T = \ell_\ttc\vect{d}$ with $\ell_\ttc = \sqrt{\norm{\vect{c}}^2 - \rho^2}$, and let $\vect{m}$ be the unit vector from $\vect{c}$ to $T$, perpendicular to $\vect{d}$. The foot of the perpendicular from $\vel_{\mathrm{rel}} = \vect{c} + \vect{w}$ is $(\vel_{\mathrm{rel}}\cdot\vect{d})\vect{d}$, and it lies beyond $T$ exactly when $\vel_{\mathrm{rel}}\cdot\vect{d} \ge \ell_\ttc = \vect{c}\cdot\vect{d}$, that is, when $\vect{w}\cdot\vect{d} \ge 0$. The front sector is bounded by the two radii to the tangent points; $\vect{w}$ outside the sector and on the side of this leg means that the angle between $\vect{w}$ and $\vect{m}$ is at most $90^\circ$ beyond the radius direction, and since $\vect{d}$ is $\vect{m}$ rotated by $90^\circ$ towards the outside of the sector, $\vect{w}\cdot\vect{d} \ge 0$ follows. With $\vect{w}$ inside the sector the arc case applies instead, so the leg formula is never used with a foot short of $T$.
\end{solution}

\begin{solution}{exr:ch13-shares}
(a) The two inequalities read $(\vel_A^{\mathrm{new}} - \vel_A - \alpha_A\vect{u})\cdot\vect{n} \ge 0$ and $(\vel_B^{\mathrm{new}} - \vel_B + \alpha_B\vect{u})\cdot(-\vect{n}) \ge 0$; their sum is $(\vel^{\mathrm{new}}_{\mathrm{rel}} - \vel_{\mathrm{rel}} - (\alpha_A + \alpha_B)\vect{u})\cdot\vect{n} \ge 0$, and the guarantee needs $(\vel^{\mathrm{new}}_{\mathrm{rel}} - \vel_{\mathrm{rel}} - \vect{u})\cdot\vect{n} \ge 0$. The difference is $(\alpha_A + \alpha_B - 1)\,\vect{u}\cdot\vect{n}$. If $\vel_{\mathrm{rel}}$ is inside the cone, $\vect{u}\cdot\vect{n} > 0$ and a sum below $1$ is not enough; if it is outside, $\vect{u}\cdot\vect{n} < 0$ and a sum above $1$ lets the pair drift too far towards the cone. Only $\alpha_A + \alpha_B = 1$ works in both cases; with $\alpha_B = 0$ the inequality of $B$ must hold with $\vel_B^{\mathrm{new}} = \vel_B$, that is, $B$ keeps the velocity that $A$ assumed. (b) A alone moves to $(1.1809, -0.1500)$, so the relative velocity is $(2.1809, -0.1500)$, at $\angle(\vel_{\mathrm{rel}}, \pos_{\mathrm{rel}}) = -11.06^\circ$ ($-3.93^\circ$ in the world frame), inside the cone of half-angle $14.36^\circ$. (c) Both take the full change: the relative velocity moves by $2\vect{u}$, twice what is needed, which is safe but is the velocity obstacle rule; at the next step each sees the other's over-reaction, the preferred velocities are permitted again, and the dance of \cref{thm:ch13-dance} returns.
\end{solution}

\begin{solution}{exr:ch13-rvo-truncated}
$2\vel - \vel_A \in \vel_B + \VO^{\ttc}_{A|B}$ holds exactly when $\vel \in \tfrac12(\vel_A + \vel_B) + \tfrac12\VO^{\ttc}_{A|B}$. By \cref{eq:ch13-vo-tau}, $\tfrac12\VO^{\ttc}_{A|B} = \bigcup_{t \in (0,\ttc]} \tfrac12\disc{\pos_{\mathrm{rel}}/t}{R/t} = \bigcup_{t \in (0,\ttc]} \disc{\pos_{\mathrm{rel}}/2t}{R/2t} = \bigcup_{t' \in (0,2\ttc]} \disc{\pos_{\mathrm{rel}}/t'}{R/t'} = \VO^{2\ttc}_{A|B}$, with $t' = 2t$. The horizon doubles because the reciprocal test looks at $2\vel - \vel_A$, a velocity twice as far from the current one, so a collision that the average velocity would cause at time $2\ttc$ shows up at $\ttc$ in the test velocity. An implementation that wants RVO agents to react only to collisions within $\ttc$ must therefore build the velocity obstacle with horizon $\ttc/2$ before applying the reciprocal transformation, or, equivalently, test $2\vel - \vel_A$ against $\VO^{\ttc/2}$; \orca avoids the issue because it works with the relative velocity directly.
\end{solution}

\begin{solution}{exr:ch13-lp-by-hand}
(a) Both orders return the same velocity, because the program has one optimum; what differs is the work. With $H_1, H_2, H_3$ the optimum for $H_1$ alone is the projection of $\vel^{\mathrm{pref}} = (1.2, 0.3)$ on the boundary of $H_1$ (it violates $H_1$), then $H_2$ is already satisfied and costs constant time, and $H_3$ is checked last. With $H_3, H_2, H_1$ the calls to the one-dimensional program come in a different order and one of them clips the interval with two earlier half-planes instead of none. The lesson is the one behind \cref{thm:ch13-lp}: the number of one-dimensional programs depends on the order, the result does not, and a random order makes the expected number small.
(b) Suppose the optimum $\vel^{\star}_{i}$ for $H_1, \dots, H_i$ lies in the interior of $H_i$ while $\vel^{\star}_{i-1}$ violates $H_i$. The feasible set is convex, so the segment from $\vel^{\star}_{i}$ to $\vel^{\star}_{i-1}$ stays inside $H_1, \dots, H_{i-1}$ and the disc, and the objective $\norm{\vel - \vel^{\mathrm{pref}}}$ decreases strictly along it (the optimum of a strictly convex objective over a convex set is unique). Moving a little way from $\vel^{\star}_{i}$ towards $\vel^{\star}_{i-1}$ stays in $H_i$ and lowers the objective, contradicting optimality. Hence $\vel^{\star}_{i}$ lies on the boundary line of $H_i$.
(c) With $\vel^{\mathrm{pref}} = (1.6, 0.4)$ the unconstrained optimum has norm $1.649 > v_{\max}$, so the first step already clamps to the disc, $\vel = 1.5\,\vel^{\mathrm{pref}}/\norm{\vel^{\mathrm{pref}}}$, and the chord test of \cref{alg:ch13-lp} decides every later one-dimensional program.
\end{solution}

\begin{solution}{exr:ch13-infeasible}
(a) At the optimum of \cref{eq:ch13-dense} the violations of the half-planes that are active are equal: if one were strictly smaller, one could move a little in the direction that decreases the largest ones without making that one the largest, and $\delta$ would drop. With three symmetric half-planes forming a triangle around the disc, the solution is therefore the centre of the disc (or the point equidistant, in violation, from the three boundary lines) and $\delta$ is that common violation.
(b) \Cref{alg:ch13-dense} walks the half-planes; when $H_i$ is violated more deeply than the current record it minimises $\mathrm{viol}_i$ subject to $\mathrm{viol}_j \le \mathrm{viol}_i$ for every earlier $j$. Each such constraint is the half-plane whose boundary is the bisector of the boundary lines of $H_i$ and $H_j$: it passes through their intersection point, where both violations vanish, with normal $\vect{n}_i - \vect{n}_j$ up to scale. Three half-planes give at most three bisectors.
(c) The objective never mentions $\vel^{\mathrm{pref}}$, so the fallback returns the same velocity for every preferred velocity: safety first. For a drone that is the right default for one or two steps but a poor steady state, because it discards the mission entirely; an implementation should keep the fallback rare (larger radii margins, a shorter $\ttc$, fewer neighbours) and hand control to the global layer of \cref{ch:ch24} when it fires repeatedly.
\end{solution}

\begin{solution}{exr:ch13-corridor}
(a) Take the nearest wall disc above the agent: $\pos_{\mathrm{rel}} = (0, 1.25)$, $R = 0.5 + 0.5 = 1$, $\norm{\pos_{\mathrm{rel}}} = 1.25$ and $\ell = \sqrt{1.25^2 - 1^2} = 0.75$. The wall does not move, so $\vel_{\mathrm{rel}} = (1,0)$; the truncating disc for $\ttc_{\mathrm{obst}} = 2$ has centre $(0, 0.625)$ and radius $0.5$, and $\vect{w} = (1, -0.625)$ has $\vect{w}\cdot\pos_{\mathrm{rel}} = -0.7813 < 0$ but $(\vect{w}\cdot\pos_{\mathrm{rel}})^2 = 0.6104 < R^2\norm{\vect{w}}^2 = 1.3906$, so \cref{thm:ch13-cases} gives a leg, and $\pos_{\mathrm{rel}} \times \vect{w} = -1.25 < 0$ makes it the right leg: $\vect{d} = (0.8, 0.6)$, $\vect{q} = (0.64, 0.48)$, $\vect{u} = (-0.36, 0.48)$, $\vect{n} = (0.6, -0.8)$. A wall takes no responsibility, so the agent takes all of $\vect{u}$ and the half-plane is $(\vel - (0.64, 0.48))\cdot(0.6, -0.8) \ge 0$. The offset vanishes because the point lies on the leg through the origin, and the constraint reduces to $0.6\,v_x - 0.8\,v_y \ge 0$, that is $v_y \le 0.75\,v_x$. The disc below gives the mirror image $v_y \ge -0.75\,v_x$, hence $\abs{v_y} \le 0.75\,v_x$: the corridor permits a lateral speed proportional to the forward speed, so an agent that has slowed to a stop can no longer step sideways.
(b) Head-on, $\vel_{\mathrm{rel}}$ points along $-\hat{\pos}$ with $\hat{\pos} = \pos_{\mathrm{rel}}/\norm{\pos_{\mathrm{rel}}}$ and the arc case applies: $\vect{u} = (d/\ttc)\hat{\pos} - \vel_{\mathrm{rel}}$ and $\vect{n} = -\hat{\pos}$, so the half-plane $(\vel - \vel_A - \tfrac12\vect{u})\cdot\vect{n} \ge 0$ reads $(\vel - \vel_A)\cdot\hat{\pos} \le \tfrac12\bigl(d/\ttc - \vel_{\mathrm{rel}}\cdot\hat{\pos}\bigr)$, which caps A's approach speed along $\hat{\pos}$ at $d/(2\ttc)$; B is capped symmetrically, so the gap closes at no more than $d/\ttc$. One step of \cref{sec:ch13-step} therefore leaves $d_{k+1} \ge d_k - \dt\,(d_k/\ttc) = d_k(1 - \dt/\ttc)$, with equality while both agents push at the cap. The factor $1 - \dt/\ttc$ is a constant in $(0,1)$, so $d_k \to 0$ geometrically and both speeds decay with it: the agents converge to the stalled state without ever reaching it, and neither ever reaches its goal.
(c) Inside the local layer, break the symmetry with the responsibility shares of \cref{sec:ch13-noncooperative}: give one agent $\alpha = 1$ and the other $\alpha = 0$ (a priority rule, or a coin flip refreshed every few seconds), so the low-priority agent takes the whole of $\vect{u}$ and backs into the bay. This resolves the pair and keeps the pairwise guarantee of \cref{thm:ch13-pairwise} because the shares still sum to one, but it guarantees nothing globally: with three agents in the corridor the priority order can cycle. In the architecture of \cref{ch:ch24}, detect the stall (speed below a threshold for several consecutive steps) and ask the global planner for a new joint plan, which is complete on the grid and can route one drone into the bay explicitly; the price is a replanning latency of seconds, which is why the local layer must keep the drones safe in the meantime.
\end{solution}
